\begin{abstract}
Performance in modern GPU-centric systems is increasingly defined by resource management policies, such as memory placement, scheduling, and observability, rather than just raw compute power. Currently, these policies are implemented at the wrong layer: either hard-coded into closed, static drivers and firmware, or siloed within user-space runtimes that lack transparency, cross-tenant visibility, and fine-grained hardware control. We argue that the OS GPU driver and device layer must serve as the system's programmable narrow waist. However, extending standard host-side eBPF is insufficient: it remains blind to critical device-side events, and naive injection into complex GPU-kernel risks system safety and efficiency.

We present \sys, a unified policy plane that treats the GPU driver and device as a single programmable OS subsystem. \sys extends GPU drivers to expose safe hooks and introduces a device-side eBPF runtime that executes verified policy logic efficiently within GPU kernels, enabling coherent, application-transparent policies that span the hardware boundary. Evaluation on realistic inference and training workloads shows that \sys improves throughput by up to  5× and reduces tail latency by 2× with negligible overhead, all without modifying applications or restarting drivers.
\end{abstract}

\section{Introduction}

GPUs dominate the compute and capital costs of modern systems, powering workloads ranging from latency-sensitive LLM inference pipelines to large-scale training, graph analytics, and vector search. Across these workloads, end-to-end performance and efficiency depend not just on raw FLOPs, but on \emph{policies}: how the system places and migrates data across device and host memory, admits and schedules kernels in multi-tenant environments, and collects observability signals for online adaptation. These policies must evolve as models, SLAs, and hardware change.

On the CPU side, programmability frameworks like eBPF already provide dynamic, safe, kernel-level policies for scheduling, caching, and observability via stable attach points\cite{Bachl21,Zhong22,jia2023programmable,schedext-docs,cacheext-docs}.
GPU stacks, however, remain dominated by vendor drivers and firmware implementing fixed, opaque mechanisms, or by profiling-focused instrumentation frameworks\cite{cupti,villa2019nvbit,neutrino,egpu}. While academic works like TimeGraph and Gdev have shifted GPU scheduling and memory control into Linux device drivers\cite{timegraph,gdev}, and recent systems further refactor driver-level control paths\cite{kernelintercept,gpreempt,lithos}, these efforts typically implement policies as static kernel modules. Lacking a safe, programmable interface, adapting policies requires invasive driver changes and operational disruptions.

To enable flexibility without kernel modification, recent serving systems retrofit programmability above the driver. Paella\cite{paella} and XSched\cite{xsched} expose software-defined scheduling, while Pie\cite{pie} and KTransformers\cite{ktransformers} offer programmable runtimes for LLM inference. While offering programmability, these user-space systems lack OS-level integration.

However, user-space approaches face three fundamental limitations. First, programmability is application-bound. Policies are tightly coupled to specific serving systems or runtimes and cannot govern arbitrary, unmodified applications. Thus, deploying a new scheduling or memory policy requires invasive code changes or porting applications to a specific stack, preventing broad reuse. Second, user-space runtimes lack global coordination for resource sharing. Confined to process silos, they cannot safely share memory or bandwidth between tenants. To avoid contention and faults, applications must conservatively pre-allocate resources (e.g., memory), preventing dynamic sharing and lowering utilization. Third, user-space systems lack fine-grained execution control. Operating above the driver, they cannot intervene at the granularity of hardware interrupts or device contexts. This confines policies to coarse-grained scheduling at kernel launch boundaries, preventing precise preemption for latency-sensitive tasks or fine-grained compute-transfer overlap.

Effective GPU resource management requires an OS-level narrow waist for policy. The GPU driver is uniquely positioned with application transparency, global cross-tenant visibility and direct control over privileged mechanisms such as page migration, kernel scheduling, and preemption. We propose treating GPU resource management as a first-class, programmable kernel subsystem by exposing a small set of verifiable policy interfaces.
Yet, applying host-only eBPF to GPUs is inadequate as critical decisions such as thread-block scheduling and warp execution progress occur device-side within GPU kernels and firmware, invisible to host-only instrumentation. A practical GPU policy substrate must separate mechanism from policy inside performance-critical drivers without destabilizing hardware interactions, and must safely execute policy logic within massively parallel, performance-sensitive GPU kernels.

Designing a general GPU policy substrate introduces three challenges. \emph{C1: Safe and efficient driver extensibility.} We must retrofit mechanism–policy separation into large, performance-critical GPU drivers without destabilizing their low-level interactions with hardware. \emph{C2: Safe and efficient device-side execution.} Policy logic must execute directly within GPU kernels on critical memory access and synchronization paths, while ensuring bounded execution and memory safety in a massively parallel environment. \emph{C3: Unified cross-layer state and management.} Effective policies require correlating host-side signals with device-side events (e.g., page faults with warp stalls), necessitating a unified state abstraction and a single control plane for verification, deployment, and debugging across the heterogeneous CPU–GPU boundary.

We present \sys, a unified, safe, and efficient policy runtime addressing these challenges. \sys exposes a small, stable set of GPU mechanisms for memory placement and migration, kernel admission and preemption, intra-kernel scheduling, and kernel-level instrumentation as programmable eBPF attach points within an extended GPU driver. \sys policies are defined as eBPF programs using a restricted instruction subset and are statically verified for memory safety and bounded execution. \sys offloads the same eBPF intermediate representation onto GPU devices: \sys compiles verified bytecode into a GPU-specific device code representation, injects it via lightweight trampolines at kernel launch time, and executes it within a GPU-side runtime. Both host-side and device-side policy programs share a common BPF map abstraction backed by unified memory, enabling low-latency, coherent state sharing without explicit data marshalling. A single verifier and control plane manage policy deployment and lifecycle across host and device, forming a unified policy domain.

We implement \sys on Linux by extending the NVIDIA open GPU kernel modules\cite{nvidia-open-gpu} to create \sysdriver and constructing a GPU-side device JIT runtime. Driver modifications leverage kernel callback interfaces to expose attach points within memory management and scheduling paths, delegating policy decisions to verified eBPF programs. On the GPU device, lightweight trampolines invoke offloaded logic with minimal overhead. Unlike prior GPU observability tools\cite{cupti,villa2019nvbit,neutrino,egpu} or user-space eBPF implementations, \sys provides a single unified, framework-agnostic policy abstraction across driver and device contexts under one verifier and control plane.

We evaluate \sys across realistic GPU workloads including LLM inference pipelines, GNN training, vector search, and mixed-priority multi-tenant scenarios. Using \sys, we implement adaptive memory prefetching and eviction policies, fine-grained kernel preemption strategies, dynamic work-stealing schedulers, and programmable GPU kernel observability tools. \sys-based policies improve throughput by up to $5\times$ and reduce tail latency by up to $2\times$ compared to static heuristics, while instrumentation-only deployments incur only $2$–$3\%$ overhead. These results confirm that unified eBPF-based programmability spanning host and device is a practical abstraction for GPU-centric systems.

\begin{itemize}[leftmargin=*,itemsep=0pt]
  \item We identify key programmability gaps in GPU stacks for memory management, scheduling, and observability, showing how the lack of a programmable, kernel-anchored policy waist at the GPU driver and device layer impedes adaptation to evolving workloads and SLAs.
  \item We design and implement \sys, an eBPF-based policy runtime treating the GPU driver and device execution layer as a unified OS subsystem. \sys exposes verified driver hooks for memory and scheduling, executes the same eBPF IR inside GPU kernels via a device-side JIT, and shares state through unified-memory-backed BPF maps under a single verifier and control plane.
  \item We demonstrate that \sys enables adaptive, workload-specific policies that significantly improve performance, resource utilization, and tail latency on realistic LLM inference, training, and vector-search workloads, with low overhead and without requiring application or driver restarts.
\end{itemize}